\usepackage{comment}



\providecommand{\plab}[1]{\textbf{#1})}
\providecommand{\pder}[2]{\frac{\partial {#1}}{\partial {#2}}}
\providecommand{\ltnorm}[1]{\left\lVert #1 \right\rVert_{2}}
\providecommand{\ghz}{\text{GHz}}

%  derivative
\providecommand{\der}[2]{\frac{\dint{#1}}{\dint{#2}}}
%  n-th derivative
\providecommand{\dern}[3]{\frac{{\dint{}^#3} #1}{{\dint{#2}^#3}}}
% partial derivative
\providecommand{\pder}[2]{\frac{\partial{#1}}{\partial{#2}}}

\providecommand{\bw}{\begin{widetext}}
        \providecommand{\ew}{\end{widetext}}

\providecommand{\newlinenn}{\\ \nn}

% poisson ratio
\providecommand{\poissonratio}{\nu_{\scalebox{0.7}{\text{P}}}}

% bold zero vector
\providecommand{\bzero}{{\bm 0}}

% quotes in italics
\providecommand{\myquote}[1]{``\textit{#1}''}


\ExplSyntaxOn
\begin{comment}
prints a collection of units of measure
Input values:
- a list of units of the parameter that appear at the numerator, e.g., \met, \second
- a list of units of the parameter that appear at the denominator (optional), e.g., \met, \second
- 'f' (optional), to be provided for the units to be displayed as a fraction

Example of usage:
\units{\pas^2, \second}[\kg, \second][f]
\end{comment}

\NewDocumentCommand{\units}{moo}
{
    \ensuremath{

        \IfNoValueTF{#3}{
            % print the units of the parameter: print the units at the numerator
            \tl_if_blank:nTF{#1}{
                % the first argument (units at the numerator ) is empty -> print `1' at the numerator
                1
            }{
                \clist_use:nn {#1} {\,}
            }
            % if the fourth argument was provided, print the units at the denominator
            \IfValueTF{#2}{
                / \int_compare:nNnTF { \clist_count:n {#2} } > { 1 }
                    { ( \clist_use:nn {#2} {\,} ) }
                    { \clist_use:nn {#2} {\,} }
            }{
            }
        }{
            % print the units of the parameter: print the units at the numerator
            \frac{\clist_use:nn {#1} {\,}}{\clist_use:nn {#2} {\,}}

        }
    }
}

\begin{comment}
returns a set of units of measure as in `units', between square braces
Input values:
- a list of units of the parameter that appear at the numerator, e.g., \met, \second
- a list of units of the parameter that appear at the denominator (optional), e.g., \met, \second
- 'f' (optional), to be provided for the units to be displayed as a fraction
Return values:
- [units]
\end{comment}

\NewDocumentCommand{\unitsbraces}{moo}
{
    \, [ \units{#1}[#2][#3] ]
}

\begin{comment}
prints  a parameter with its units, where units at the numerator and at the denominator are separated with a slash: this method takes three arguments, the parameter value, the units at the numerator and the units at the denominator

Input values:
- the mantissa of the parameter (optional): if not provided, the mantissa is assumed to be `1'
- the exponent of 10^ if the parameter has to be printed in scientific notation (optional)
- a list of units of the parameter that appear at the numerator, e.g., \met, \second
- a list of units of the parameter that appear at the denominator (optional), e.g., \met, \second
- 'f' (optional), to be provided for the units to be displayed as a fraction

Examples of usage:

\parameter[1.3]{\met^2, \second}
\parameter[1.3]{\met^2, \second}[\pas, \second][f]
\parameter[1.3][-2]{\met, \second}[\pas, \second]
\end{comment}


\NewDocumentCommand{\parameter}{oomoo}
{
    \ensuremath{
        % print the mantissa of the parameter value

        \tl_if_blank:nTF {#1}{
            % no mantissa was provided -> print `1' as mantissa



            % if an exponent was provided, print 10^{#2}
            \IfValueTF{#2}{
                10^{\clist_map_inline:nn {#2} {##1}}
            }{
                1
            }



        }{
            % a mantissa was provided -> print it


            % if an exponent was provided, print 10^{#2}
            \IfValueTF{#2}{
                \clist_map_inline:nn {#1} {##1}  \times 10^{\clist_map_inline:nn {#2} {##1}}
            }{
                \clist_map_inline:nn {#1} {##1}
            }

        }
        % print the units of the parameter: print the units at the numerator
        \, \units{#3}[#4][#5]
    }
}
\ExplSyntaxOff


